% !TeX root = ../main.tex

\renewcommand\nomgroup[1]{%
  \item[\bfseries
  \ifstrequal{#1}{A}{Stichprobe und deskriptive Statistik}{%
  \ifstrequal{#1}{B}{Inferenzstatistik}{%
  }}%
]}

% Stichprobenumfang: N für die gesamte Stichprobe, n für eine Gruppe oder Analyse.
\nomenclature[Aa]{$N$}{Gesamtumfang der Stichprobe}
\nomenclature[Ab]{$n$}{Anzahl der gültigen Beobachtungen in der jeweiligen Gruppe oder Analyse}

% Lage und Streuung der erhobenen Werte; M und SD beziehen sich auf die Stichprobe.
\nomenclature[Ac]{$M$}{Arithmetischer Mittelwert der Stichprobe}
\nomenclature[Ad]{$SD$}{Standardabweichung der Stichprobe}
\nomenclature[Ae]{$Mdn$}{Median}
\nomenclature[Af]{$IQR$}{Interquartilsabstand: Differenz zwischen dem dritten und dem ersten Quartil}

% Signifikanzniveau und p-Wert für die spätere Prüfung statistischer Hypothesen.
\nomenclature[Ba]{$\alpha$}{Festgelegtes Signifikanzniveau}
\nomenclature[Bb]{$p$}{p-Wert eines statistischen Tests}

% Konfidenzintervalle beschreiben die Unsicherheit einer Schätzung.
\nomenclature[Bc]{KI}{Konfidenzintervall; das Konfidenzniveau wird bei der jeweiligen Angabe genannt}
